\documentclass{article}
\usepackage[utf8]{inputenc}

\title{\textbf{Speaker Recognition or Identification using Machine Learning}

\underline{\Large{4-week report}}}
\author{\textbf{Adhiraj Banerjee(SRFP Application No. : ENGS589)}}
\date{July 2021}

\begin{document}

\maketitle

\section{\textbf{\underline{Introduction}:}}
The topic for my Internship which was assigned to me by my guide is, \textit{Speaker Recognition or Identification using Machine Learning}. Here, the problem I have to solve is that I have to create a trained model which will be able to recognize the speaker from its voice sample provided to it as its input. Speaker Recognition is important in today's world. It is used in various applications, like, voice biometric, forensic applictions, etc. Speaker Recognition determines which registered speaker is providing its speech as input among the set of all registered speakers. In order to solve this problem, I have considered the \textbf{Machine Learning Approach} as a possible solution for it.

\section{\textbf{\underline{The Machine Learning Approach As a Solution}:}}
 \textit{Machine Learning} is a method which helps systems to learn from the training data so that they can predict the class of the test data given as input, without being explicitly programmed for this purpose. This same method would be applied in the Speaker Recognition System. Using Machine Learning, a model for Speaker Recognition will be trained in such a way that it will be able to identify unique patterns and features from different voice samples and distinguish a voice sample from the rest. This method can be applied by, firstly \textbf{gathering different people's speech samples}, \textbf{extract features from the voice sample} suitable for the classifier, which is \textbf{trained for building a trained model} and finally,\textbf{performing classification for recognition of each Speaker} from its corresponding input voice sample.

\section{\textbf{\underline{Conclusion}:}}
I would like to conclude this report by stating that, the \textit{Machine Learning Approach} is such a practical solution, which makes \textbf{Speaker Recognition} more advantageous over \textbf{Speech Recognition}. Speech Recognition is a \textbf{text-dependent} method, where it highly depends on the language and corpus, whereas, Speaker Recognition mainly focusses on the raw audio's percepts and its related data by which it identifies the uniqueness among different speakers. Hence, \textbf{Machine Learning} helps us to achieve the goal of Speaker Recognition and makes us understand about differences between \textit{Speaker Recognition} \& \textit{Speech Recognition}.


\end{document}